% Part: first-order-logic
% Chapter: proof-systems
% Section: natural-deduction

\documentclass[../../../include/open-logic-section]{subfiles}

\begin{document}

\iftag{FOL}
      {\olfileid{fol}{prf}{ntd}}
      {\olfileid{pl}{prf}{ntd}}

\olsection{الاستنتاج الطبيعي}

الاستنتاج الطبيعي نظام !!a{derivation} يُقصد به محاكاة الاستدلال
الفعلي (ولا سيما نمط الاستدلال المنضبط الذي يستعمله الرياضيون).
فالاستدلال الفعلي يسير وفق عدد من الأنماط «الطبيعية». فمثلًا، يتيح لنا
البرهان بالحالات إثبات نتيجة انطلاقًا من مقدمة فصلية، وذلك بإثبات أن
النتيجة تتبع من كل واحد من طرفي الفصل. ويتيح لنا البرهان غير المباشر
إثبات نتيجة بإظهار أن نفيها يؤدي إلى تناقض. ويثبت البرهان الشرطي دعوى
شرطية «إذا \dots فإن \dots» بإظهار أن التالي يتبع من المقدَّم.
والاستنتاج الطبيعي تصيير صوري لبعض هذه الاستدلالات الطبيعية. فلكل من
الروابط المنطقية والمكممات قاعدتان، قاعدة إدخال وقاعدة حذف، وتقابل كل
منهما نمطًا طبيعيًا من هذا القبيل. فمثلًا، تقابل~$\Intro{\lif}$
البرهان الشرطي، وتقابل~$\Elim{\lor}$ البرهان بالحالات. ومن القواعد
البسيطة بوجه خاص القاعدة~$\Elim{\land}$، التي تجيز الاستدلال من
$!A \land !B$ إلى~$!A$ (أو~$!B$).

ومن السمات التي تميز الاستنتاج الطبيعي من سائر أنظمة ال!!{derivation}
استعماله للفروض. و!!^a{derivation} في الاستنتاج الطبيعي شجرة من
ال!!{formula}s. وتقع !!{formula} واحدة في جذر شجرة ال!!{formula}s،
أما «أوراق» الشجرة فهي ال!!{formula}s التي تُشتق منها النتيجة. وفي
الاستنتاج الطبيعي تؤدي بعض !!{formula}s الواقعة في الأوراق دورًا داخل
ال!!{derivation}، لكن دورها يكون قد «استُنفد» حين يصل
ال!!{derivation} إلى النتيجة. وهذا يقابل الممارسة الفعلية المتمثلة في
إدخال فرضيات لا تبقى نافذة إلا مدة قصيرة. ففي البرهان بالحالات،
مثلًا، نفترض صدق كل من طرفي الفصل؛ وفي البرهان الشرطي نفترض صدق
المقدَّم؛ وفي البرهان غير المباشر نفترض صدق نفي النتيجة. وهذه الطريقة
في إدخال فروض افتراضية ثم الاستغناء عنها من أجل إثبات خطوة وسيطة سمة
مميزة للاستنتاج الطبيعي. وتسمى الصيغ الواقعة في أوراق
!!{derivation} بالاستنتاج الطبيعي فروضًا، وقد تجيز بعض قواعد
الاستدلال «!!{discharge}» هذه الفروض. فمثلًا، إذا كان لدينا
!!a{derivation} لـ~$!B$ من فروض تتضمن~$!A$، فإن القاعدة
$\Intro{\lif}$ تجيز لنا استنتاج~$!A \lif !B$ وإسقاط أي فرض
من الصورة~$!A$. (ولكي نتتبع الفروض التي يجري إسقاطها عند كل استدلال،
نوسم الاستدلال والفروض التي يسقطها بعدد.) وتكفي الفروض التي يستمر
!!{undischarged}ها في نهاية ال!!{derivation}، مجتمعةً، لصدق النتيجة؛
ومن ثم يثبت !!a{derivation} أن فروضه التي يستمر !!{undischarged}ها
تستلزم نتيجته.

تتحقق العلاقة~$\Gamma \Proves !A$ القائمة على الاستنتاج الطبيعي إذا
وفقط إذا وجد !!a{derivation} تكون فيه~$!A$ آخر !!{sentence} في
الشجرة، وكان كل ورقة يستمر !!{undischarged}ها عنصرًا في~$\Gamma$.
وتكون~$!A$ مبرهنة في الاستنتاج الطبيعي إذا وفقط إذا وجد
!!a{derivation} تكون فيه~$!A$ آخر !!{sentence}، وقد جرى
!!{discharged} جميع الفروض. فمثلًا، هذا !!a{derivation} يبين أن
$\Proves (!A \land !B) \lif !A$:
\begin{prooftree}
  \AxiomC{$\Discharge{!A \land !B}{1}$}
  \RightLabel{\Elim{\land}}
  \UnaryInfC{$!A$}
  \DischargeRule{\Intro{\lif}}{1}
  \UnaryInfC{$(!A \land !B) \lif !A$}
\end{prooftree}
يبين الوسم~$1$ أن الفرض~$!A \land !B$ قد جرى !!{discharged}ه عند
استدلال~\Intro{\lif}.
  
تكون المجموعة~$\Gamma$ غير متسقة إذا وفقط إذا كان
$\Gamma \Proves \lfalse$ في الاستنتاج الطبيعي. وتجعل القاعدة
\FalseInt{} أي !!{sentence} قابلةً لل!!{derive} من مجموعة غير
متسقة.

طوّر غيرهارد غنتسن وستانيسواف ياشكوفسكي أنظمة الاستنتاج الطبيعي في
ثلاثينيات القرن العشرين، ثم طورها لاحقًا داغ برافيتس وفريدريك فيتش.
ولأن استدلالاتها تحاكي طرق البرهان الطبيعية، فإن الفلاسفة يؤثرونها.
وكثيرًا ما تُستعمل النسخ التي طورها فيتش في الكتب التمهيدية في
المنطق. وقد عُدّت قواعد الاستنتاج الطبيعي أحيانًا، في فلسفة المنطق،
محددةً لمعاني المؤثرات المنطقية («الدلالة القائمة على نظرية
البرهان»).

\end{document}
